\begin{figure}
	\resizebox{\textwidth}{!}{
		\centering
		\subcaptionbox{%Marges de phase pour une raideur nominale
			Raideur nominale ($K_n=194$)\label{subfig_theory_stability_map_nominal_stiffness}}
		{
			\tikzsetnextfilename{subfig_theory_stability_map_nominal}	
			\begin{tikzpicture} 
				\begin{axis}[%
					width=.45\linewidth,
					height=.43\linewidth,
					xlabel={$K_p$ (\si{\radian\per\second\per\newton})},
					ylabel={$K_i$ (\si{\radian\squared\per\second\squared\per\newton})},
					xlabel style={sloped},
					ylabel style={sloped},
					%zlabel={$\Delta\phi$ (\si{\degree})},
					zlabel={$\Delta_M$},
					%zticklabel={\pgfmathparse{\tick*100}\pgfmathprintnumber{\pgfmathresult}},
					colormap name=ShirazToOrient,
					scale only axis,  
					view = {0}{90},
					%point meta max=90,
					%point meta min=-20,
					point meta max=1,
					point meta min=0,
					%ztick={0,0.25,0.5,0.75,1},
					clip=false,
					scaled x ticks={base 10:2},
					]
					%
					%\addplot3 [surf, mesh/rows=40, shader=interp, patch type=bilinear] table[col sep=comma, x=ki, y=kp, z=Dphi] {misc/chap_5/stability_map_nominal_theory.csv};
					\addplot3 [surf, mesh/rows=50, shader=interp, patch type=bilinear] table[col sep=comma, y=ki, x=kp, z=zeta] {misc/chap_5/damping_map_nominal_theory.csv};
					\node at (axis cs:0.015,0.08,0.52) {\pgfuseplotmark{*}};
					\draw[dashed, black] (axis cs:0,0) rectangle (axis cs:2.2e-2,0.3) node[below, xshift=-0.5cm] {exp.};
				\end{axis}
			\end{tikzpicture}
		}
		\subcaptionbox{%Marges de phase pour une raideur multipliée par 3
			Raideur supposée en cocontraction ($K_r=3K_n=582$)\label{subfig_theory_stability_map_high_stiffness}}
		{
			\tikzsetnextfilename{subfig_theory_stability_map_cocontraction}
			\begin{tikzpicture} 
				\begin{axis}[%
					name = ax2,
					width=.45\linewidth,
					height=.43\linewidth,
					xlabel={$K_p$ (\si{\radian\per\second\per\newton})},
					%ylabel={$Kp$},
					xlabel style={sloped},
					ylabel style={sloped},
					%zlabel={$\Delta\phi$ (\si{\degree})},
					zlabel={$\Delta_M$},
					%xticklabel={\pgfmathparse{\tick*100}\pgfmathprintnumber{\pgfmathresult}},
					colormap name=ShirazToOrient,
					scale only axis,  
					view = {0}{90},
					%xmax=299,
					%ymin=250,
					%point meta max=90,
					%point meta min=-20,
					point meta max=1,
					point meta min=0,
					%ztick={0,45,90,135},
					%ztick={0,0.25,0.5,0.75,1},
					yticklabels={,,},
					clip=false,
					scaled x ticks={base 10:2},
					]
					%
					%\addplot3 [surf, mesh/rows=40, shader=interp, patch type=bilinear] table[col sep=comma, x=ki, y=kp, z=Dphi] {misc/chap_5/stability_map_stiff_theory.csv};
					\addplot3 [surf, mesh/rows=50, shader=interp, patch type=bilinear] table[col sep=comma, y=ki, x=kp, z=zeta] {misc/chap_5/damping_map_stiff_theory.csv};
					\node at (axis cs:0.015,0.08,0.86) {\pgfuseplotmark{*}};
					\draw[dashed, black] (axis cs:0,0) rectangle (axis cs:2.2e-2,0.3);
					\draw[dotted, white, thick] (axis cs:0,0.3) -- (axis cs:0.009,0.8); -- (axis cs:0,0.8) -- (axis cs:0,0.3);
					%\node at (rel axis cs:0,0.5,1) [above] {\Large{$K$}};
				\end{axis}
				\begin{axis}[
					name=ax3,
					hide axis,
					scale only axis,
					width=0pt,
					height=0pt,
					colormap name=ShirazToOrient,
					colorbar horizontal,					
					point meta max=1,
					point meta min=0,
					colorbar style={
						width=6.5cm,
						rotate=90,
						%title= $\Delta_M$,
						at={($(ax2.south east) + (0.1cm,0)$)},
						anchor=south west, 
						%xtick style={draw=none},
						xticklabel style={xshift=.4cm},
					},
					]
					\addplot [draw=none] coordinates {(0,0) (1,1)};
				\end{axis}				
				\node at ($(ax2.south east) + (0.8cm,-0.5cm)$) {$\zeta$};
			\end{tikzpicture}
		}
	}
	\caption[Reproduction théorique de la cartographie expérimentale de stabilité du robot contrôlé en admittance]{Reproduction théorique de la cartographie expérimentale (zone entourée en tirets noirs) utilisant les paramètres du participant 000 de la phase \cref{item_pert2}, en fonction des gains du correcteur cartésien en admittance. La palette de couleurs indique le coefficient d'amortissement $\zeta$ le plus faible des pôles du système en boucle fermée \cref{eq_close_loop_system}. La zone délimitée par des pointillés blancs est instable.}
	\label{fig_theory_stability_map}
\end{figure}